\subsection{Sliding Window}
The sliding window approach is a method for processing steaming data using a fixed -size buffer that retains only the most recent observations. As new data arrives, older data points are discarded while more recent data points are continuously appended to the  buffer. \\ In the context of this thesis, the sliding window approach is applied to the time series data before the Time delay Embedding (Section 2.2.1) and Principal Component Analysis (Section 2.2.2). This fulfills the key requirement of handling steaming data efficiently while maintaining constant memory usage  and low latency.  \\
According (\cite{https://doi.org/10.1155/2014/879736} p.14 ), given a time series $T$ of length $m$ and a user-defined subsequence length $n$, all possible subsequences $C_p$ can be extracted by sliding a window of size $n$ across $T$.
This formalization reflects how sliding window mechanism extracts sequential segments from continuous time-series data.